
\documentclass{article}
\usepackage{colortbl}
\usepackage{makecell}
\usepackage{multirow}
\usepackage{supertabular}

\begin{document}

\newcounter{utterance}

\centering \large Interaction Transcript for game `cladder', experiment `full\_v1.5\_default', episode 8915 with qwen3.5{-}9b.
\vspace{24pt}

{ \footnotesize  \setcounter{utterance}{1}
\setlength{\tabcolsep}{0pt}
\begin{supertabular}{c@{$\;$}|p{.15\linewidth}@{}p{.15\linewidth}p{.15\linewidth}p{.15\linewidth}p{.15\linewidth}p{.15\linewidth}}
    \# & \multicolumn{2}{c}{Player} && \multicolumn{2}{c}{Game Master} \\
    \hline

    \theutterance \stepcounter{utterance}  
    & & & \multicolumn{4}{p{0.6\linewidth}}{
        \cellcolor[rgb]{0.9,0.9,0.9}{
            \makecell[{{p{\linewidth}}}]{
                \texttt{\tiny{[P1$\langle$GM]}}
                \texttt{You are an expert in causal inference. The following question is not a typical commonsense query, but rather a meticulously designed question created by a professor specializing in causal inference, intended to assess the students' mastery of the course content.} \\
\texttt{Imagine a self{-}contained, hypothetical world with only the following conditions, and without any unmentioned factors or causal relationships: Zuph has a direct effect on wibl. Uvzi has a direct effect on wibl. Method 1: We look directly at how zuph correlates with uvzi in general. Method 2: We look at this correlation case by case according to wibl. To understand how zuph affects uvzi, is it more correct to use the Method 1 than Method 2?} \\
\texttt{Start your answer with 'Yes' or 'No', followed by additional reasoning or evidence to support your explanation.} \\
            }
        }
    }
     \\ \\

    \theutterance \stepcounter{utterance}  
    & \multicolumn{4}{p{0.6\linewidth}}{
        \cellcolor[rgb]{0.561,0.737,0.561}{
            \makecell[{{p{\linewidth}}}]{
                \texttt{\tiny{[P1$\rangle$GM]}}
                \texttt{No.} \\
\\ 
\texttt{In this hypothetical world, the causal structure is defined as follows: \$Zuph \textbackslash{}rightarrow Wibl\$ and \$Uvzi \textbackslash{}rightarrow Wibl\$. There are no arrows connecting \$Zuph\$ and \$Uvzi\$ directly, nor is there any common cause mentioned between them. Therefore, in the absence of any specified causal link, the correlation between \$Zuph\$ and \$Uvzi\$ is entirely spurious (non{-}causal).} \\
\\ 
\texttt{**Method 1** looks at the marginal correlation between \$Zuph\$ and \$Uvzi\$. In this specific setup, unless there is an unmentioned confounder or a specific functional dependency creating a statistical artifact, this correlation should be zero (or random noise). Relying on this to understand how \$Zuph\$ affects \$Uvzi\$ is incorrect because it ignores the fact that they are only connected via the common effect \$Wibl\$.} \\
\\ 
\texttt{**Method 2** looks at the correlation between \$Zuph\$ and \$Uvzi\$ conditional on \$Wibl\$ (i.e., stratifying by \$Wibl\$). This is the correct approach to assess the relationship between two variables that share a common effect. By conditioning on the common effect \$Wibl\$, we block the "collider path" (\$Zuph \textbackslash{}rightarrow Wibl \textbackslash{}leftarrow Uvzi\$). In causal inference, conditioning on a collider induces a spurious association between its parents if one is not already present, or it removes the spurious association induced by the collider if one exists. More importantly, since there is no direct causal path between \$Zuph\$ and \$Uvzi\$, the conditional correlation (after blocking the collider) should reveal the true lack of causal influence, whereas the marginal correlation might be misleadingly non{-}zero due to the collider bias.} \\
\\ 
\texttt{Therefore, to correctly determine if \$Zuph\$ affects \$Uvzi\$ (which it does not), one must account for the structure involving \$Wibl\$. Method 2 is the rigorous application of causal adjustment (specifically, controlling for a collider to avoid bias or to isolate the direct effect, which is zero here), whereas Method 1 risks inferring a causal relationship from a spurious correlation induced by the common effect. Thus, Method 2 is more correct.} \\
            }
        }
    }
    & & \\ \\

    \theutterance \stepcounter{utterance}  
    & & & \multicolumn{2}{p{0.3\linewidth}}{
        \cellcolor[rgb]{0.9,0.9,0.9}{
            \makecell[{{p{\linewidth}}}]{
                \texttt{\tiny{[GM$|$GM]}}
                \texttt{No} \\
            }
        }
    }
    & & \\ \\

    \theutterance \stepcounter{utterance}  
    & & & \multicolumn{2}{p{0.3\linewidth}}{
        \cellcolor[rgb]{0.9,0.9,0.9}{
            \makecell[{{p{\linewidth}}}]{
                \texttt{\tiny{[GM$|$GM]}}
                \texttt{yes} \\
            }
        }
    }
    & & \\ \\

    \theutterance \stepcounter{utterance}  
    & & & \multicolumn{2}{p{0.3\linewidth}}{
        \cellcolor[rgb]{0.9,0.9,0.9}{
            \makecell[{{p{\linewidth}}}]{
                \texttt{\tiny{[GM$|$GM]}}
                \texttt{game\_result = LOSE} \\
            }
        }
    }
    & & \\ \\

\end{supertabular}
}

\end{document}
